% =============================================================================
% UNIFIED STYLE GUIDE
% =============================================================================

\newpage
\section{Style Guide and Conventions}
\label{sec:style-guide}

This section consolidates coding conventions, naming standards, and best practices for all STAR project components. Following these guidelines ensures consistency across firmware, protocol definitions, and documentation.

\subsection{Terminology Standard}

The project uses inclusive terminology following OSHWA (Open Source Hardware Association) standards.

\begin{table}[H]
\centering
\caption{Communication Bus Terminology}
\label{tab:inclusive-terminology}
\begin{tabular}{|l|l|p{6cm}|}
\hline
\textbf{Use} & \textbf{Avoid} & \textbf{Context} \\
\hline
Controller/Peripheral & Master/Slave & I2C, SPI, 1-Wire bus roles \\
COPI & MOSI & Controller Out, Peripheral In \\
CIPO & MISO & Controller In, Peripheral Out \\
Primary/Main & Master & Configuration structures \\
\hline
\end{tabular}
\end{table}

\layerbox{keyinsightcolor}{ESP-IDF Legacy API Mapping}{%
ESP-IDF APIs still use legacy terminology internally (e.g., \texttt{I2C\_MODE\_SLAVE}, \texttt{mosi\_io\_num}). Map these to OSHWA terminology in comments and documentation. Never propagate legacy terms into new code or documentation.
}

\subsection{Naming Conventions}

\subsubsection{C/C++ Firmware (ESP32)}

\begin{table}[H]
\centering
\caption{ESP32 Firmware Naming Conventions}
\label{tab:esp32-naming}
\begin{tabular}{|l|l|l|}
\hline
\textbf{Element} & \textbf{Convention} & \textbf{Example} \\
\hline
Functions & snake\_case & \texttt{star\_motor\_set\_duty()} \\
Variables & snake\_case & \texttt{velocity\_rpm}, \texttt{current\_ma} \\
Macros/Constants & SCREAMING\_SNAKE & \texttt{MAX\_MOTOR\_CURRENT\_MA} \\
Types & snake\_case\_t suffix & \texttt{star\_motor\_handle\_t} \\
Static functions & internal\_ prefix & \texttt{internal\_update\_pid()} \\
Private functions & priv\_ prefix & \texttt{priv\_calculate\_crc()} \\
Static variables & s\_ prefix & \texttt{s\_motor\_count} \\
Global variables & g\_ prefix (avoid) & \texttt{g\_system\_state} \\
\hline
\end{tabular}
\end{table}

\subsubsection{Protocol Buffers (Boston Dynamics Style)}

\begin{table}[H]
\centering
\caption{Protocol Buffer Naming Conventions}
\label{tab:proto-naming-style}
\begin{tabular}{|l|l|l|}
\hline
\textbf{Element} & \textbf{Convention} & \textbf{Example} \\
\hline
Package & Versioned namespace & \texttt{star.v1} \\
Messages & PascalCase & \texttt{VelocityCommand}, \texttt{BatteryState} \\
Fields & snake\_case + unit suffix & \texttt{velocity\_mps}, \texttt{temperature\_celsius} \\
Enums & SCREAMING\_SNAKE\_CASE & \texttt{MOTOR\_STATE\_RUNNING} \\
Enum Zero Value & \_UNKNOWN suffix & \texttt{STATUS\_UNKNOWN = 0} \\
Enum Values & Type name prefix & \texttt{ROBOT\_MODE\_MANUAL} \\
RPC Pattern & \{Verb\}Request/Response & \texttt{SetVelocityRequest} \\
\hline
\end{tabular}
\end{table}

\subsection{Unit Suffixes (MKS System)}

All numeric fields use SI units with explicit suffixes. This eliminates unit ambiguity and enables automatic documentation generation.

\begin{table}[H]
\centering
\caption{Standard Unit Suffixes}
\label{tab:unit-suffixes-style}
\begin{tabular}{|l|l|l|l|}
\hline
\textbf{Suffix} & \textbf{Unit} & \textbf{Proto Example} & \textbf{C Example} \\
\hline
\texttt{\_m} & meters & \texttt{x\_m} & \texttt{distance\_m} \\
\texttt{\_mps} & meters/second & \texttt{velocity\_mps} & \texttt{speed\_mps} \\
\texttt{\_mps2} & m/s$^2$ & \texttt{accel\_x\_mps2} & \texttt{acceleration\_mps2} \\
\texttt{\_rad} & radians & \texttt{angle\_rad} & \texttt{heading\_rad} \\
\texttt{\_rad\_per\_s} & rad/second & \texttt{omega\_rad\_per\_s} & \texttt{angular\_vel\_rad\_per\_s} \\
\texttt{\_celsius} & degrees C & \texttt{temperature\_celsius} & \texttt{motor\_temp\_celsius} \\
\texttt{\_deci\_celsius} & 0.1\textdegree C & \texttt{temp\_deci\_celsius} & (BMS native format) \\
\texttt{\_us} & microseconds & \texttt{timestamp\_us} & \texttt{elapsed\_us} \\
\texttt{\_ms} & milliseconds & \texttt{timeout\_ms} & \texttt{period\_ms} \\
\texttt{\_ma} & milliamps & \texttt{current\_ma} & \texttt{motor\_current\_ma} \\
\texttt{\_mv} & millivolts & \texttt{cell\_mv} & \texttt{pack\_voltage\_mv} \\
\texttt{\_mw} & milliwatts & \texttt{power\_mw} & \texttt{consumption\_mw} \\
\texttt{\_percent} & 0--100\% & \texttt{soc\_percent} & \texttt{duty\_cycle\_percent} \\
\hline
\end{tabular}
\end{table}

\subsection{ESP32 Firmware Architecture}

The firmware follows \textbf{SOLID principles} with Dependency Inversion (DIP) as the foundation.

\subsubsection{Dependency Inversion Pattern}

\begin{enumerate}[leftmargin=2em]
    \item \textbf{Core Abstraction Layer} (\texttt{star\_core/}) -- Interfaces only, no implementations
    \item \textbf{Implementation Modules} -- Concrete implementations of interfaces
    \item \textbf{Injection Pattern} -- Components receive interfaces, not concrete types
\end{enumerate}

\layerbox{layer1color}{DIP Benefits for Embedded Systems}{%
\begin{itemize}[leftmargin=1.5em,topsep=2pt,itemsep=1pt]
    \item Components can be tested with mock interfaces (no hardware required)
    \item Implementations can be swapped without changing dependent code
    \item Optional dependencies can be passed as NULL
    \item No circular dependencies between modules
\end{itemize}
}

\subsubsection{Memory Management Rules}

\begin{itemize}[leftmargin=1.5em]
    \item \textbf{No dynamic allocation} -- All buffers statically allocated
    \item \textbf{Avoid large stack allocations} -- Use static/global for large arrays
    \item \textbf{Use \texttt{const}} -- Put lookup tables and defaults in flash/ROM
    \item \textbf{Validate all buffers} -- Check lengths before \texttt{memcpy}, DMA, ring buffer ops
\end{itemize}

\subsubsection{Critical Coding Rules}

\begin{enumerate}[leftmargin=2em]
    \item \textbf{Always use braces} for control statements (if, for, while), even single statements
    \item \textbf{Assertions for programming errors only} -- Use \texttt{ESP\_ERROR\_CHECK} for ESP-IDF, proper error handling for runtime errors
    \item \textbf{Avoid inline ASM} -- If required, use \texttt{volatile} and document why
    \item \textbf{Minimize ISR work} -- Set flags or queue work to background tasks
    \item \textbf{No blocking delays} -- Use event-driven design or non-blocking drivers
\end{enumerate}

\subsection{Embedded Systems Best Practices}

\subsubsection{Real-Time and Safety}

\begin{table}[H]
\centering
\caption{Real-Time Best Practices}
\label{tab:realtime-practices}
\begin{tabular}{|l|p{9cm}|}
\hline
\textbf{Practice} & \textbf{Implementation} \\
\hline
Know timing constraints & Profile worst-case execution times (WCET) \\
Use watchdog timers & Automatic reset on firmware lockup \\
Handle sensor timeouts & Never assume sensors always respond \\
Fail safe & On error, disable actuators and enter safe state \\
CRC critical data & Validate packets and stored configuration \\
\hline
\end{tabular}
\end{table}

\subsubsection{Hardware Abstraction}

\begin{itemize}[leftmargin=1.5em]
    \item Use abstract bus interfaces (\texttt{star\_bus}) for I2C, SPI, UART, GPIO, ADC
    \item Debounce mechanical inputs with time-based filtering
    \item Enable brown-out detection (BOD/BOR)
    \item Implement flash wear-leveling for frequently written data
    \item Always use pull-up/pull-down on GPIO inputs (no floating pins)
\end{itemize}

\subsection{Protocol Buffer Constraints (nanopb)}

ESP32 requires static allocation for all protocol buffer fields. Configure sizes in \texttt{.options} files.

\begin{table}[H]
\centering
\caption{nanopb Field Size Constraints}
\label{tab:nanopb-constraints}
\begin{tabular}{|l|l|l|}
\hline
\textbf{Field Type} & \textbf{Constraint} & \textbf{Example} \\
\hline
String (UUID) & \texttt{max\_size:64} & \texttt{request\_id} \\
String (error) & \texttt{max\_size:128} & \texttt{error\_message} \\
String (version) & \texttt{max\_size:32} & \texttt{firmware\_version} \\
Repeated (cells) & \texttt{max\_count:16} & \texttt{cell\_mv[]} \\
Repeated (motors) & \texttt{max\_count:4} & \texttt{motor\_status[]} \\
Bytes (firmware) & \texttt{max\_size:1024} & \texttt{chunk\_data} \\
\hline
\end{tabular}
\end{table}

\subsection{Documentation Standards}

\begin{itemize}[leftmargin=1.5em]
    \item Use Doxygen-compatible comments for all public APIs
    \item Document the ``why'' not just the ``what''
    \item Keep configuration centralized in header files
    \item Avoid magic numbers -- use named constants with comments
    \item Version hardware and firmware together using git tags
\end{itemize}

